\begin{table}[htbp]
\centering
\begin{threeparttable}
\begin{tabular}{lccc}
\toprule
Outcome & 60 s (N=163) & 150 s (N=72) & 300 s (N=45) \\
\midrule
Mean HR & 15.994 [14.774, 17.214] & 17.685 [15.976, 19.395] & 19.891 [17.105, 22.677] \\
SDNN & -10.742 [-11.871, -9.613] & -11.024 [-12.734, -9.314] & -12.267 [-14.680, -9.854] \\
RMSSD & -0.034 [-0.089, 0.022] & -0.054 [-0.147, 0.040] & -0.104 [-0.244, 0.036] \\
LF power & -1.395 [-1.524, -1.267] & -1.537 [-1.742, -1.332] & -1.722 [-2.034, -1.410] \\
HF power & -0.789 [-0.947, -0.631] & -0.924 [-1.185, -0.664] & -1.109 [-1.486, -0.733] \\
LF/HF & -0.603 [-0.723, -0.483] & -0.606 [-0.796, -0.416] & -0.600 [-0.807, -0.393] \\
Sample Entropy & 0.368 [0.292, 0.443] & 0.320 [0.182, 0.458] & 0.316 [0.134, 0.498] \\
\bottomrule
\end{tabular}\begin{tablenotes}\footnotesize\item Entries are VO$_2$ coefficients per one SD increase with 95\% CI. Thinning was a sensitivity analysis; the complete AR(1)-GEE used all 923 observations.\end{tablenotes}
\end{threeparttable}
\caption{Temporal-thinning sensitivity}
\label{tab:s1-thinning}
\end{table}
